% ===== Фінал =====
\begin{darkframe}
\begin{frame}[plain,t]
\begin{slidecanvas}
  \kerbstrip{(-0.3,-9)}{34}{0.28}
  \checker{(12.4,-3.0)}{9}{7}{0.4}
  \fill[dark,path fading=west] (12.3,-3.1) rectangle ++(1.4,3.2);
  \node[anchor=north west,align=left,text width=9cm] at (0.9,-1.6) {%
    {\fontsize{40}{40}\selectfont\heavy\color{white}Дякуємо}\par\vskip0.3cm
    {\large\light\color{white!80!dark}Питання?}};
  \node[anchor=north west,align=left,font=\scriptsize,text=later,text width=11cm] at (0.9,-4.3) {%
    {\semibold\color{kerb}\caps{Артефакти етапу 3}}\par\vskip0.15cm
    {\ttfamily docs/stages/stage-3/}\par\vskip0.08cm
    {\color{white}Project-Overview.md}\enspace опис, цілі, обсяг, обмеження, ризики\\
    {\color{white}Roadmap.md, gantt.pdf}\enspace перелік робіт і діаграма Ґанта\\
    {\color{white}RACI.md}\enspace розподіл відповідальності і навантаження\\
    {\color{white}Backlog.pdf}\enspace беклог продукту\\
    {\color{white}Eisenhower.pdf}\enspace матриця Ейзенхауера};
\end{slidecanvas}
\note{\speaker{Лід · Яна}
Це наш план на рік. Усі артефакти лежать у репозиторії, у папці docs/stages/stage-3. Дякуємо за увагу — будемо раді питанням.}
\end{frame}
\end{darkframe}
